%%%%%%%%%%%%%%%%%%%%%%%%%%%%%%%%%%%%%%%%%%%%%%%%%%%%%%%%%%%%
% 5.14 DISTRIBUTION THEORY
% The Composition of Advanced Mathematics
%%%%%%%%%%%%%%%%%%%%%%%%%%%%%%%%%%%%%%%%%%%%%%%%%%%%%%%%%%%%

\section{Distribution Theory}
\label{sec:distribution-theory}

\prerequisite{
Real analysis, measure theory, Lebesgue integration, functional analysis,
multivariable calculus, basic topology, and differential equations.
}

\learningobjectives{
By the end of this section, the reader should be able to:
\begin{itemize}
    \item explain why generalized functions are needed;
    \item work with smooth compactly supported test functions;
    \item define distributions as continuous linear functionals;
    \item identify locally integrable functions with regular distributions;
    \item distinguish regular and singular distributions;
    \item work with the Dirac delta distribution;
    \item define distributional derivatives;
    \item compute derivatives of discontinuous functions;
    \item understand weak derivatives;
    \item multiply distributions by smooth functions;
    \item translate and scale distributions;
    \item understand support and singular support;
    \item define convergence of distributions;
    \item use tensor products and distributions in several variables;
    \item understand convolution involving distributions;
    \item identify fundamental solutions of differential operators;
    \item interpret differential equations distributionally;
    \item understand the connection with Fourier transforms;
    \item recognize tempered distributions;
    \item work with the Schwartz space;
    \item understand principal-value distributions;
    \item recognize the relationship between distributions and Sobolev spaces;
    \item apply distribution theory to PDEs, physics, Fourier analysis,
          and generalized differential equations.
\end{itemize}
}

%%%%%%%%%%%%%%%%%%%%%%%%%%%%%%%%%%%%%%%%%%%%%%%%%%%%%%%%%%%%
\subsection{Why Generalized Functions?}
\label{subsec:why-generalized-functions}
%%%%%%%%%%%%%%%%%%%%%%%%%%%%%%%%%%%%%%%%%%%%%%%%%%%%%%%%%%%%

Classical analysis studies ordinary functions such as

\[
f:\R\to\R.
\]

However, many important objects appearing in differential equations and
physics do not behave like ordinary functions.

The most famous example is the Dirac delta.

Informally, one would like an object satisfying

\[
\boxed{
\delta(x)=0
\quad\text{for }x\neq0,
}
\]

while also satisfying

\[
\boxed{
\int_{-\infty}^{\infty}
\delta(x)\,dx
=
1.
}
\]

No ordinary function has these properties in the intended classical
sense.

Distribution theory replaces the idea of evaluating such an object
pointwise with the idea of studying how it acts on test functions.

Thus,

\[
\boxed{
\text{distribution}
=
\text{generalized function acting on test functions}.
}
\]

%%%%%%%%%%%%%%%%%%%%%%%%%%%%%%%%%%%%%%%%%%%%%%%%%%%%%%%%%%%%
\subsection{Test Functions}
\label{subsec:distribution-test-functions}
%%%%%%%%%%%%%%%%%%%%%%%%%%%%%%%%%%%%%%%%%%%%%%%%%%%%%%%%%%%%

Let

\[
\Omega\subseteq\R^n
\]

be open.

\begin{definitionbox}{Test Function}
A test function on \(\Omega\) is an infinitely differentiable function

\[
\varphi:\Omega\to\mathbb{R}
\]

or

\[
\varphi:\Omega\to\mathbb{C}
\]

whose support is compact and contained in \(\Omega\).
\end{definitionbox}

The lowercase Greek letter \(\varphi\), read \emph{phi}, is commonly used
for test functions.

The space of test functions is denoted

\[
\boxed{
C_c^\infty(\Omega).
}
\]

The notation means

\[
\boxed{
C^\infty
=
\text{infinitely differentiable},
}
\]

and

\[
\boxed{
c
=
\text{compact support}.
}
\]

Another common notation is

\[
\boxed{
\mathcal{D}(\Omega)
=
C_c^\infty(\Omega).
}
\]

%%%%%%%%%%%%%%%%%%%%%%%%%%%%%%%%%%%%%%%%%%%%%%%%%%%%%%%%%%%%
\subsection{Support of a Function}
\label{subsec:function-support-distributions}
%%%%%%%%%%%%%%%%%%%%%%%%%%%%%%%%%%%%%%%%%%%%%%%%%%%%%%%%%%%%

\begin{definitionbox}{Support}
The support of a function \(f\) is

\[
\boxed{
\operatorname{supp}(f)
=
\overline{
\{x:f(x)\neq0\}
}.
}
\]
\end{definitionbox}

Thus a compactly supported function is zero outside some compact subset.

Test functions therefore vanish outside a bounded region strictly inside
the domain.

This property eliminates boundary contributions in many integration-by-parts
arguments.

%%%%%%%%%%%%%%%%%%%%%%%%%%%%%%%%%%%%%%%%%%%%%%%%%%%%%%%%%%%%
\subsection{The Test-Function Space}
\label{subsec:test-function-space}
%%%%%%%%%%%%%%%%%%%%%%%%%%%%%%%%%%%%%%%%%%%%%%%%%%%%%%%%%%%%

The space

\[
\mathcal{D}(\Omega)
\]

is more than an ordinary vector space.

Its topology encodes both

\[
\boxed{
\text{uniform convergence of every derivative}
}
\]

and

\[
\boxed{
\text{control of supports}.
}
\]

For a fixed compact set \(K\subseteq\Omega\), a sequence
\(\varphi_j\) converges to \(0\) in the corresponding test-function
space when all \(\varphi_j\) are supported in \(K\) and

\[
\boxed{
\sup_{x\in K}
|\partial^\alpha\varphi_j(x)|
\to0
}
\]

for every multi-index \(\alpha\).

This topology is essential to the definition of a distribution.

%%%%%%%%%%%%%%%%%%%%%%%%%%%%%%%%%%%%%%%%%%%%%%%%%%%%%%%%%%%%
\subsection{Multi-Index Notation}
\label{subsec:distribution-multi-index}
%%%%%%%%%%%%%%%%%%%%%%%%%%%%%%%%%%%%%%%%%%%%%%%%%%%%%%%%%%%%

For functions of several variables, let

\[
\alpha
=
(\alpha_1,\ldots,\alpha_n)
\]

where each \(\alpha_j\) is a nonnegative integer.

Define

\[
\boxed{
|\alpha|
=
\alpha_1+\cdots+\alpha_n.
}
\]

The derivative corresponding to \(\alpha\) is

\[
\boxed{
\partial^\alpha
=
\frac{
\partial^{|\alpha|}
}{
\partial x_1^{\alpha_1}
\cdots
\partial x_n^{\alpha_n}
}.
}
\]

Multi-index notation greatly simplifies formulas in distribution theory
and partial differential equations.

%%%%%%%%%%%%%%%%%%%%%%%%%%%%%%%%%%%%%%%%%%%%%%%%%%%%%%%%%%%%
\subsection{Distributions}
\label{subsec:definition-distributions}
%%%%%%%%%%%%%%%%%%%%%%%%%%%%%%%%%%%%%%%%%%%%%%%%%%%%%%%%%%%%

\begin{definitionbox}{Distribution}
A distribution on an open set \(\Omega\subseteq\R^n\) is a continuous
linear functional

\[
\boxed{
T:C_c^\infty(\Omega)\to\mathbb{F},
}
\]

where

\[
\mathbb{F}\in\{\R,\mathbb{C}\}.
\]
\end{definitionbox}

The value of \(T\) on a test function \(\varphi\) is commonly written

\[
\boxed{
\langle T,\varphi\rangle.
}
\]

This is read

\[
\boxed{
\text{``the distribution \(T\) acting on \(\varphi\).''}
}
\]

The space of distributions is denoted

\[
\boxed{
\mathcal{D}'(\Omega).
}
\]

The prime indicates the continuous dual of the test-function space.

%%%%%%%%%%%%%%%%%%%%%%%%%%%%%%%%%%%%%%%%%%%%%%%%%%%%%%%%%%%%
\subsection{Linearity of Distributions}
\label{subsec:distribution-linearity}
%%%%%%%%%%%%%%%%%%%%%%%%%%%%%%%%%%%%%%%%%%%%%%%%%%%%%%%%%%%%

A distribution satisfies

\[
\boxed{
\langle T,
a\varphi+b\psi\rangle
=
a\langle T,\varphi\rangle
+
b\langle T,\psi\rangle.
}
\]

Thus distributions are linear objects.

This linear structure allows differential equations to be extended from
ordinary functions to generalized functions.

%%%%%%%%%%%%%%%%%%%%%%%%%%%%%%%%%%%%%%%%%%%%%%%%%%%%%%%%%%%%
\subsection{Regular Distributions}
\label{subsec:regular-distributions}
%%%%%%%%%%%%%%%%%%%%%%%%%%%%%%%%%%%%%%%%%%%%%%%%%%%%%%%%%%%%

Ordinary functions can define distributions.

Suppose

\[
f\in L_{\mathrm{loc}}^1(\Omega).
\]

The notation

\[
L_{\mathrm{loc}}^1(\Omega)
\]

means that \(f\) is integrable on every compact subset of \(\Omega\).

Define

\[
\boxed{
\langle T_f,\varphi\rangle
=
\int_\Omega
f(x)\varphi(x)\,dx.
}
\]

Then \(T_f\) is a distribution.

Such distributions are called regular distributions.

%%%%%%%%%%%%%%%%%%%%%%%%%%%%%%%%%%%%%%%%%%%%%%%%%%%%%%%%%%%%
\subsection{Locally Integrable Functions}
\label{subsec:locally-integrable-functions}
%%%%%%%%%%%%%%%%%%%%%%%%%%%%%%%%%%%%%%%%%%%%%%%%%%%%%%%%%%%%

\begin{definitionbox}{Locally Integrable Function}
A measurable function \(f\) belongs to

\[
L_{\mathrm{loc}}^1(\Omega)
\]

if

\[
\boxed{
\int_K|f(x)|\,dx<\infty
}
\]

for every compact set

\[
K\subseteq\Omega.
\]
\end{definitionbox}

Local integrability is enough because every test function has compact
support.

Therefore

\[
\int_\Omega f\varphi\,dx
\]

is well-defined.

%%%%%%%%%%%%%%%%%%%%%%%%%%%%%%%%%%%%%%%%%%%%%%%%%%%%%%%%%%%%
\subsection{Functions Equal Almost Everywhere}
\label{subsec:distributions-ae}
%%%%%%%%%%%%%%%%%%%%%%%%%%%%%%%%%%%%%%%%%%%%%%%%%%%%%%%%%%%%

If

\[
f=g
\quad\text{almost everywhere},
\]

then

\[
\boxed{
T_f=T_g.
}
\]

Indeed,

\[
\int_\Omega
f\varphi\,dx
=
\int_\Omega
g\varphi\,dx
\]

for every test function \(\varphi\).

Thus regular distributions naturally identify functions that agree
almost everywhere.

%%%%%%%%%%%%%%%%%%%%%%%%%%%%%%%%%%%%%%%%%%%%%%%%%%%%%%%%%%%%
\subsection{Singular Distributions}
\label{subsec:singular-distributions}
%%%%%%%%%%%%%%%%%%%%%%%%%%%%%%%%%%%%%%%%%%%%%%%%%%%%%%%%%%%%

Not every distribution arises from a locally integrable function.

Such distributions are called singular distributions.

The standard example is the Dirac delta distribution.

%%%%%%%%%%%%%%%%%%%%%%%%%%%%%%%%%%%%%%%%%%%%%%%%%%%%%%%%%%%%
\subsection{Dirac Delta Distribution}
\label{subsec:dirac-delta}
%%%%%%%%%%%%%%%%%%%%%%%%%%%%%%%%%%%%%%%%%%%%%%%%%%%%%%%%%%%%

\begin{definitionbox}{Dirac Delta Distribution}
The Dirac delta at \(a\in\R^n\), denoted \(\delta_a\), is defined by

\[
\boxed{
\langle\delta_a,\varphi\rangle
=
\varphi(a).
}
\]
\end{definitionbox}

At the origin one commonly writes

\[
\boxed{
\delta
=
\delta_0.
}
\]

Thus

\[
\boxed{
\langle\delta,\varphi\rangle
=
\varphi(0).
}
\]

The delta distribution extracts the value of a test function at one
point.

%%%%%%%%%%%%%%%%%%%%%%%%%%%%%%%%%%%%%%%%%%%%%%%%%%%%%%%%%%%%
\subsection{The Delta Is Not an Ordinary Function}
\label{subsec:delta-not-function}
%%%%%%%%%%%%%%%%%%%%%%%%%%%%%%%%%%%%%%%%%%%%%%%%%%%%%%%%%%%%

The notation

\[
\delta(x)
\]

is useful symbolically, but it should not be interpreted as an ordinary
function with an infinitely large value at \(0\).

The rigorous meaning is entirely contained in

\[
\boxed{
\langle\delta,\varphi\rangle
=
\varphi(0).
}
\]

Similarly, the informal identity

\[
\int
\delta(x)\varphi(x)\,dx
=
\varphi(0)
\]

is shorthand for the distributional action.

%%%%%%%%%%%%%%%%%%%%%%%%%%%%%%%%%%%%%%%%%%%%%%%%%%%%%%%%%%%%
\subsection{Shifted Delta}
\label{subsec:shifted-delta}
%%%%%%%%%%%%%%%%%%%%%%%%%%%%%%%%%%%%%%%%%%%%%%%%%%%%%%%%%%%%

The delta concentrated at \(a\) satisfies

\[
\boxed{
\langle\delta_a,\varphi\rangle
=
\varphi(a).
}
\]

Symbolically,

\[
\boxed{
\delta_a
=
\delta(x-a).
}
\]

Thus the delta can represent a quantity concentrated at a single point.

%%%%%%%%%%%%%%%%%%%%%%%%%%%%%%%%%%%%%%%%%%%%%%%%%%%%%%%%%%%%
\subsection{Measures as Distributions}
\label{subsec:measures-as-distributions}
%%%%%%%%%%%%%%%%%%%%%%%%%%%%%%%%%%%%%%%%%%%%%%%%%%%%%%%%%%%%

A locally finite measure \(\mu\) can define a distribution by

\[
\boxed{
\langle T_\mu,\varphi\rangle
=
\int_\Omega
\varphi\,d\mu.
}
\]

For example, the Dirac measure at \(a\) satisfies

\[
\int\varphi\,d\delta_a
=
\varphi(a).
\]

Thus measure theory and distribution theory overlap naturally.

%%%%%%%%%%%%%%%%%%%%%%%%%%%%%%%%%%%%%%%%%%%%%%%%%%%%%%%%%%%%
\subsection{Distributional Derivatives}
\label{subsec:distributional-derivatives}
%%%%%%%%%%%%%%%%%%%%%%%%%%%%%%%%%%%%%%%%%%%%%%%%%%%%%%%%%%%%

The most important feature of distributions is that every distribution
can be differentiated.

Suppose first that \(f\) is smooth and compactly supported.

Integration by parts gives

\[
\int
f'(x)\varphi(x)\,dx
=
-
\int
f(x)\varphi'(x)\,dx.
\]

This motivates the general definition.

\begin{definitionbox}{Distributional Derivative}
The derivative \(T'\) of a distribution \(T\) is defined by

\[
\boxed{
\langle T',\varphi\rangle
=
-
\langle T,\varphi'\rangle.
}
\]
\end{definitionbox}

%%%%%%%%%%%%%%%%%%%%%%%%%%%%%%%%%%%%%%%%%%%%%%%%%%%%%%%%%%%%
\subsection{Higher Distributional Derivatives}
\label{subsec:higher-distributional-derivatives}
%%%%%%%%%%%%%%%%%%%%%%%%%%%%%%%%%%%%%%%%%%%%%%%%%%%%%%%%%%%%

Repeated differentiation gives

\[
\boxed{
\langle T^{(m)},\varphi\rangle
=
(-1)^m
\langle T,\varphi^{(m)}\rangle.
}
\]

In several variables,

\[
\boxed{
\langle
\partial^\alpha T,
\varphi
\rangle
=
(-1)^{|\alpha|}
\langle
T,
\partial^\alpha\varphi
\rangle.
}
\]

Therefore every distribution possesses derivatives of every order.

%%%%%%%%%%%%%%%%%%%%%%%%%%%%%%%%%%%%%%%%%%%%%%%%%%%%%%%%%%%%
\subsection{Why Distributional Differentiation Works}
\label{subsec:why-distributional-differentiation}
%%%%%%%%%%%%%%%%%%%%%%%%%%%%%%%%%%%%%%%%%%%%%%%%%%%%%%%%%%%%

Ordinary differentiation can fail when a function has corners, jumps,
or other singularities.

Distributional differentiation transfers derivatives onto smooth test
functions.

Since test functions are infinitely differentiable,

\[
\boxed{
\text{the derivative acts on the test function instead}.
}
\]

This allows differentiation to continue beyond the classical setting.

%%%%%%%%%%%%%%%%%%%%%%%%%%%%%%%%%%%%%%%%%%%%%%%%%%%%%%%%%%%%
\subsection{Derivative of the Dirac Delta}
\label{subsec:delta-derivative}
%%%%%%%%%%%%%%%%%%%%%%%%%%%%%%%%%%%%%%%%%%%%%%%%%%%%%%%%%%%%

By definition,

\[
\langle\delta',\varphi\rangle
=
-
\langle\delta,\varphi'\rangle.
\]

Therefore,

\[
\boxed{
\langle\delta',\varphi\rangle
=
-\varphi'(0).
}
\]

More generally,

\[
\boxed{
\langle\delta^{(m)},\varphi\rangle
=
(-1)^m
\varphi^{(m)}(0).
}
\]

%%%%%%%%%%%%%%%%%%%%%%%%%%%%%%%%%%%%%%%%%%%%%%%%%%%%%%%%%%%%
\subsection{Heaviside Function}
\label{subsec:heaviside-function}
%%%%%%%%%%%%%%%%%%%%%%%%%%%%%%%%%%%%%%%%%%%%%%%%%%%%%%%%%%%%

Define the Heaviside step function by

\[
H(x)
=
\begin{cases}
0, & x<0,\\
1, & x>0.
\end{cases}
\]

The value at \(x=0\) does not affect the corresponding regular
distribution.

Classically, \(H'\) is zero wherever \(x\neq0\), while the derivative
does not exist at \(0\).

Distributionally, however,

\[
\boxed{
H'=\delta.
}
\]

%%%%%%%%%%%%%%%%%%%%%%%%%%%%%%%%%%%%%%%%%%%%%%%%%%%%%%%%%%%%
\subsection{Worked Example: Derivative of the Heaviside Function}
\label{subsec:worked-heaviside-derivative}
%%%%%%%%%%%%%%%%%%%%%%%%%%%%%%%%%%%%%%%%%%%%%%%%%%%%%%%%%%%%

\begin{workedexamplebox}{Distributional Derivative of a Step Function}

For a test function \(\varphi\),

\[
\langle H',\varphi\rangle
=
-
\langle H,\varphi'\rangle.
\]

Since \(H=1\) on \((0,\infty)\),

\[
\langle H',\varphi\rangle
=
-
\int_0^\infty
\varphi'(x)\,dx.
\]

Because \(\varphi\) has compact support,

\[
\varphi(x)=0
\]

for sufficiently large \(x\).

Therefore,

\[
-
\int_0^\infty
\varphi'(x)\,dx
=
\varphi(0).
\]

Hence

\begin{answerbox}
\[
\boxed{
H'=\delta.
}
\]
\end{answerbox}

\end{workedexamplebox}

%%%%%%%%%%%%%%%%%%%%%%%%%%%%%%%%%%%%%%%%%%%%%%%%%%%%%%%%%%%%
\subsection{Derivative of a Function with a Jump}
\label{subsec:jump-derivative}
%%%%%%%%%%%%%%%%%%%%%%%%%%%%%%%%%%%%%%%%%%%%%%%%%%%%%%%%%%%%

Suppose \(f\) is piecewise smooth and has a jump at \(x=a\).

Let

\[
[f]_a
=
f(a^+)-f(a^-).
\]

Then its distributional derivative contains a delta contribution:

\[
\boxed{
D f
=
f'_{\mathrm{classical}}
+
[f]_a\delta_a,
}
\]

where the classical derivative is interpreted separately on the smooth
pieces.

Thus jumps become delta distributions after differentiation.

%%%%%%%%%%%%%%%%%%%%%%%%%%%%%%%%%%%%%%%%%%%%%%%%%%%%%%%%%%%%
\subsection{Absolute Value Example}
\label{subsec:absolute-value-distribution}
%%%%%%%%%%%%%%%%%%%%%%%%%%%%%%%%%%%%%%%%%%%%%%%%%%%%%%%%%%%%

Consider

\[
f(x)=|x|.
\]

Classically,

\[
f'(x)
=
\begin{cases}
-1, & x<0,\\
1, & x>0.
\end{cases}
\]

Thus

\[
f'
=
\operatorname{sgn}(x)
\]

as a distribution.

Differentiating again,

\[
\boxed{
f''=2\delta.
}
\]

Therefore a corner in a function produces a delta after two
distributional derivatives.

%%%%%%%%%%%%%%%%%%%%%%%%%%%%%%%%%%%%%%%%%%%%%%%%%%%%%%%%%%%%
\subsection{Weak Derivatives}
\label{subsec:weak-derivatives}
%%%%%%%%%%%%%%%%%%%%%%%%%%%%%%%%%%%%%%%%%%%%%%%%%%%%%%%%%%%%

Distributional differentiation leads directly to weak derivatives.

\begin{definitionbox}{Weak Derivative}
Let \(f,g\in L_{\mathrm{loc}}^1(\Omega)\).

We say \(g\) is the weak derivative of \(f\) with respect to \(x_j\) if

\[
\boxed{
\int_\Omega
f\,\partial_j\varphi\,dx
=
-
\int_\Omega
g\varphi\,dx
}
\]

for every

\[
\varphi\in C_c^\infty(\Omega).
\]
\end{definitionbox}

Thus \(g\) is a weak derivative precisely when the distributional
derivative of \(f\) is represented by the locally integrable function
\(g\).

%%%%%%%%%%%%%%%%%%%%%%%%%%%%%%%%%%%%%%%%%%%%%%%%%%%%%%%%%%%%
\subsection{Classical and Weak Derivatives}
\label{subsec:classical-weak-derivatives}
%%%%%%%%%%%%%%%%%%%%%%%%%%%%%%%%%%%%%%%%%%%%%%%%%%%%%%%%%%%%

If \(f\) is continuously differentiable, then its weak derivative agrees
with its classical derivative.

Thus,

\[
\boxed{
\text{classical derivative}
\Longrightarrow
\text{weak derivative}.
}
\]

The converse need not hold.

Weak derivatives allow functions with much lower regularity to be
differentiated.

%%%%%%%%%%%%%%%%%%%%%%%%%%%%%%%%%%%%%%%%%%%%%%%%%%%%%%%%%%%%
\subsection{Connection with Sobolev Spaces}
\label{subsec:distribution-sobolev}
%%%%%%%%%%%%%%%%%%%%%%%%%%%%%%%%%%%%%%%%%%%%%%%%%%%%%%%%%%%%

Sobolev spaces are defined using weak derivatives.

For example,

\[
\boxed{
W^{1,p}(\Omega)
}
\]

consists of functions

\[
u\in L^p(\Omega)
\]

whose first weak derivatives also belong to \(L^p(\Omega)\).

The Hilbert-space case

\[
\boxed{
H^1(\Omega)
=
W^{1,2}(\Omega)
}
\]

is especially important in variational methods and PDEs.

Distribution theory therefore provides the derivative concept underlying
Sobolev spaces.

%%%%%%%%%%%%%%%%%%%%%%%%%%%%%%%%%%%%%%%%%%%%%%%%%%%%%%%%%%%%
\subsection{Multiplication by Smooth Functions}
\label{subsec:multiplying-distributions}
%%%%%%%%%%%%%%%%%%%%%%%%%%%%%%%%%%%%%%%%%%%%%%%%%%%%%%%%%%%%

A distribution can be multiplied by a smooth function.

If

\[
a\in C^\infty(\Omega)
\]

and

\[
T\in\mathcal{D}'(\Omega),
\]

define

\[
\boxed{
\langle aT,\varphi\rangle
=
\langle T,a\varphi\rangle.
}
\]

For the delta distribution,

\[
\boxed{
a\delta_a
=
a(a)\delta_a.
}
\]

%%%%%%%%%%%%%%%%%%%%%%%%%%%%%%%%%%%%%%%%%%%%%%%%%%%%%%%%%%%%
\subsection{Product Rule}
\label{subsec:distribution-product-rule}
%%%%%%%%%%%%%%%%%%%%%%%%%%%%%%%%%%%%%%%%%%%%%%%%%%%%%%%%%%%%

If \(a\) is smooth and \(T\) is a distribution, then

\[
\boxed{
D(aT)
=
a'T+aT'.
}
\]

More generally, multi-index differentiation satisfies the corresponding
Leibniz rule.

Thus many familiar differentiation identities extend naturally when one
factor is smooth.

%%%%%%%%%%%%%%%%%%%%%%%%%%%%%%%%%%%%%%%%%%%%%%%%%%%%%%%%%%%%
\subsection{Products of Arbitrary Distributions}
\label{subsec:products-arbitrary-distributions}
%%%%%%%%%%%%%%%%%%%%%%%%%%%%%%%%%%%%%%%%%%%%%%%%%%%%%%%%%%%%

There is no general canonical multiplication operation

\[
\boxed{
\mathcal{D}'\times\mathcal{D}'
\to
\mathcal{D}'.
}
\]

In particular, expressions such as

\[
\delta^2
\]

are not defined in ordinary distribution theory without additional
structure.

This is a major limitation of the theory.

%%%%%%%%%%%%%%%%%%%%%%%%%%%%%%%%%%%%%%%%%%%%%%%%%%%%%%%%%%%%
\subsection{Translation of Distributions}
\label{subsec:translation-distributions}
%%%%%%%%%%%%%%%%%%%%%%%%%%%%%%%%%%%%%%%%%%%%%%%%%%%%%%%%%%%%

For a function,

\[
f(x-a)
\]

represents translation by \(a\).

Distributional translation is defined through test functions.

If \(\tau_aT\) denotes translation, then

\[
\boxed{
\langle\tau_aT,\varphi\rangle
=
\langle T,\varphi(\,\cdot+a)\rangle.
}
\]

With this convention,

\[
\boxed{
\tau_a\delta_0
=
\delta_a.
}
\]

%%%%%%%%%%%%%%%%%%%%%%%%%%%%%%%%%%%%%%%%%%%%%%%%%%%%%%%%%%%%
\subsection{Scaling the Delta Distribution}
\label{subsec:delta-scaling}
%%%%%%%%%%%%%%%%%%%%%%%%%%%%%%%%%%%%%%%%%%%%%%%%%%%%%%%%%%%%

For a nonzero scalar \(a\),

\[
\boxed{
\delta(ax)
=
\frac{1}{|a|}
\delta(x)
}
\]

in one dimension.

More generally, if

\[
A:\R^n\to\R^n
\]

is invertible, then

\[
\boxed{
\delta(Ax)
=
\frac{1}{|\det A|}
\delta(x).
}
\]

These identities are understood distributionally.

%%%%%%%%%%%%%%%%%%%%%%%%%%%%%%%%%%%%%%%%%%%%%%%%%%%%%%%%%%%%
\subsection{Delta of a Function}
\label{subsec:delta-of-function}
%%%%%%%%%%%%%%%%%%%%%%%%%%%%%%%%%%%%%%%%%%%%%%%%%%%%%%%%%%%%

Suppose \(g:\R\to\R\) is smooth and has simple zeros

\[
x_1,\ldots,x_m,
\]

meaning

\[
g(x_j)=0,
\qquad
g'(x_j)\neq0.
\]

Then formally and distributionally,

\[
\boxed{
\delta(g(x))
=
\sum_{j=1}^{m}
\frac{\delta(x-x_j)}
{|g'(x_j)|}.
}
\]

This identity appears frequently in physics and change-of-variable
calculations.

%%%%%%%%%%%%%%%%%%%%%%%%%%%%%%%%%%%%%%%%%%%%%%%%%%%%%%%%%%%%
\subsection{Support of a Distribution}
\label{subsec:support-distribution}
%%%%%%%%%%%%%%%%%%%%%%%%%%%%%%%%%%%%%%%%%%%%%%%%%%%%%%%%%%%%

A distribution \(T\) vanishes on an open set \(U\subseteq\Omega\) if

\[
\boxed{
\langle T,\varphi\rangle=0
}
\]

for every

\[
\varphi\in C_c^\infty(U).
\]

The support of \(T\) is the complement of the largest open set on which
\(T\) vanishes.

For example,

\[
\boxed{
\operatorname{supp}(\delta_a)=\{a\}.
}
\]

%%%%%%%%%%%%%%%%%%%%%%%%%%%%%%%%%%%%%%%%%%%%%%%%%%%%%%%%%%%%
\subsection{Singular Support}
\label{subsec:singular-support}
%%%%%%%%%%%%%%%%%%%%%%%%%%%%%%%%%%%%%%%%%%%%%%%%%%%%%%%%%%%%

The singular support of a distribution identifies where it fails to be
represented locally by a smooth function.

It is denoted

\[
\boxed{
\operatorname{sing\,supp}(T).
}
\]

For the Dirac delta,

\[
\boxed{
\operatorname{sing\,supp}(\delta_0)
=
\{0\}.
}
\]

For the Heaviside distribution,

\[
\boxed{
\operatorname{sing\,supp}(H)
=
\{0\}.
}
\]

Singular support is an elementary precursor to the more refined concept
of wavefront sets.

%%%%%%%%%%%%%%%%%%%%%%%%%%%%%%%%%%%%%%%%%%%%%%%%%%%%%%%%%%%%
\subsection{Convergence of Distributions}
\label{subsec:distribution-convergence}
%%%%%%%%%%%%%%%%%%%%%%%%%%%%%%%%%%%%%%%%%%%%%%%%%%%%%%%%%%%%

A sequence of distributions \(T_n\) converges to \(T\) if

\[
\boxed{
\langle T_n,\varphi\rangle
\to
\langle T,\varphi\rangle
}
\]

for every test function

\[
\varphi\in C_c^\infty(\Omega).
\]

This is called convergence in the sense of distributions.

It is a weak form of convergence because it is defined through testing
against smooth functions.

%%%%%%%%%%%%%%%%%%%%%%%%%%%%%%%%%%%%%%%%%%%%%%%%%%%%%%%%%%%%
\subsection{Approximation of the Delta}
\label{subsec:approximation-delta}
%%%%%%%%%%%%%%%%%%%%%%%%%%%%%%%%%%%%%%%%%%%%%%%%%%%%%%%%%%%%

Let

\[
\rho\in C_c^\infty(\R^n)
\]

satisfy

\[
\rho\geq0
\]

and

\[
\int_{\R^n}\rho(x)\,dx=1.
\]

Define

\[
\boxed{
\rho_\varepsilon(x)
=
\frac{1}{\varepsilon^n}
\rho\left(\frac{x}{\varepsilon}\right).
}
\]

Then

\[
\boxed{
\rho_\varepsilon
\to
\delta
}
\]

in the sense of distributions as

\[
\varepsilon\to0^+.
\]

%%%%%%%%%%%%%%%%%%%%%%%%%%%%%%%%%%%%%%%%%%%%%%%%%%%%%%%%%%%%
\subsection{Worked Example: Delta Approximation}
\label{subsec:worked-delta-approximation}
%%%%%%%%%%%%%%%%%%%%%%%%%%%%%%%%%%%%%%%%%%%%%%%%%%%%%%%%%%%%

\begin{workedexamplebox}{Approximate Identity Converging to Delta}

For a test function \(\varphi\),

\[
\langle\rho_\varepsilon,\varphi\rangle
=
\int_{\R^n}
\frac{1}{\varepsilon^n}
\rho\left(\frac{x}{\varepsilon}\right)
\varphi(x)\,dx.
\]

Set

\[
x=\varepsilon y.
\]

Then

\[
dx=\varepsilon^n\,dy,
\]

so

\[
\langle\rho_\varepsilon,\varphi\rangle
=
\int_{\R^n}
\rho(y)\varphi(\varepsilon y)\,dy.
\]

As

\[
\varepsilon\to0,
\]

we have

\[
\varphi(\varepsilon y)\to\varphi(0).
\]

Therefore,

\[
\langle\rho_\varepsilon,\varphi\rangle
\to
\varphi(0)
\int_{\R^n}\rho(y)\,dy.
\]

Since the integral of \(\rho\) is \(1\),

\begin{answerbox}
\[
\boxed{
\rho_\varepsilon
\to
\delta
\quad
\text{in }\mathcal{D}'.
}
\]
\end{answerbox}

\end{workedexamplebox}

%%%%%%%%%%%%%%%%%%%%%%%%%%%%%%%%%%%%%%%%%%%%%%%%%%%%%%%%%%%%
\subsection{Mollifiers}
\label{subsec:mollifiers}
%%%%%%%%%%%%%%%%%%%%%%%%%%%%%%%%%%%%%%%%%%%%%%%%%%%%%%%%%%%%

Functions such as \(\rho_\varepsilon\) are called mollifiers.

If \(f\) is locally integrable, define

\[
\boxed{
f_\varepsilon
=
f*\rho_\varepsilon.
}
\]

Under standard hypotheses,

\[
f_\varepsilon
\]

is smooth.

As

\[
\varepsilon\to0,
\]

the smoothed functions approximate \(f\) in appropriate senses.

Mollifiers are fundamental tools for approximation and regularization.

%%%%%%%%%%%%%%%%%%%%%%%%%%%%%%%%%%%%%%%%%%%%%%%%%%%%%%%%%%%%
\subsection{Convolution of Functions}
\label{subsec:distribution-convolution-functions}
%%%%%%%%%%%%%%%%%%%%%%%%%%%%%%%%%%%%%%%%%%%%%%%%%%%%%%%%%%%%

For suitable functions \(f\) and \(g\), their convolution is

\[
\boxed{
(f*g)(x)
=
\int_{\R^n}
f(x-y)g(y)\,dy.
}
\]

Convolution combines information from nearby points and often produces a
smoother function.

For the delta distribution,

\[
\boxed{
f*\delta=f.
}
\]

Thus the delta acts as the identity element for convolution.

%%%%%%%%%%%%%%%%%%%%%%%%%%%%%%%%%%%%%%%%%%%%%%%%%%%%%%%%%%%%
\subsection{Convolution with Distributions}
\label{subsec:convolution-distributions}
%%%%%%%%%%%%%%%%%%%%%%%%%%%%%%%%%%%%%%%%%%%%%%%%%%%%%%%%%%%%

Convolution can be extended to distributions under appropriate support
conditions.

For example, if \(T\) is a distribution and \(\varphi\) is a test
function, then \(T*\varphi\) is a smooth function.

A useful schematic formula is

\[
\boxed{
(T*\varphi)(x)
=
\langle
T,
\varphi(x-\cdot)
\rangle.
}
\]

Convolution therefore provides a mechanism for smoothing distributions.

%%%%%%%%%%%%%%%%%%%%%%%%%%%%%%%%%%%%%%%%%%%%%%%%%%%%%%%%%%%%
\subsection{Differentiation and Convolution}
\label{subsec:differentiation-convolution-distributions}
%%%%%%%%%%%%%%%%%%%%%%%%%%%%%%%%%%%%%%%%%%%%%%%%%%%%%%%%%%%%

Under appropriate hypotheses,

\[
\boxed{
\partial^\alpha(T*\varphi)
=
(\partial^\alpha T)*\varphi
=
T*(\partial^\alpha\varphi).
}
\]

This identity allows derivatives to be transferred between factors.

It is fundamental in PDE theory.

%%%%%%%%%%%%%%%%%%%%%%%%%%%%%%%%%%%%%%%%%%%%%%%%%%%%%%%%%%%%
\subsection{Differential Operators}
\label{subsec:distribution-differential-operators}
%%%%%%%%%%%%%%%%%%%%%%%%%%%%%%%%%%%%%%%%%%%%%%%%%%%%%%%%%%%%

A linear differential operator has the form

\[
\boxed{
P(x,D)
=
\sum_{|\alpha|\leq m}
a_\alpha(x)\partial^\alpha.
}
\]

If the coefficients \(a_\alpha\) are smooth, the operator acts naturally
on distributions.

Thus

\[
\boxed{
P(x,D)T
}
\]

is defined distributionally even when \(T\) is not an ordinary
differentiable function.

%%%%%%%%%%%%%%%%%%%%%%%%%%%%%%%%%%%%%%%%%%%%%%%%%%%%%%%%%%%%
\subsection{Constant-Coefficient Differential Operators}
\label{subsec:constant-coefficient-operators}
%%%%%%%%%%%%%%%%%%%%%%%%%%%%%%%%%%%%%%%%%%%%%%%%%%%%%%%%%%%%

A constant-coefficient differential operator has the form

\[
\boxed{
P(D)
=
\sum_{|\alpha|\leq m}
c_\alpha\partial^\alpha.
}
\]

Examples include

\[
\boxed{
\frac{d}{dx},
}
\]

\[
\boxed{
\Delta
=
\sum_{j=1}^{n}
\frac{\partial^2}{\partial x_j^2},
}
\]

and the wave operator

\[
\boxed{
\frac{\partial^2}{\partial t^2}
-
c^2\Delta.
}
\]

Distribution theory provides a natural setting for these operators.

%%%%%%%%%%%%%%%%%%%%%%%%%%%%%%%%%%%%%%%%%%%%%%%%%%%%%%%%%%%%
\subsection{Fundamental Solutions}
\label{subsec:fundamental-solutions}
%%%%%%%%%%%%%%%%%%%%%%%%%%%%%%%%%%%%%%%%%%%%%%%%%%%%%%%%%%%%

\begin{definitionbox}{Fundamental Solution}
A distribution \(E\) is a fundamental solution for a differential
operator \(P(D)\) if

\[
\boxed{
P(D)E=\delta.
}
\]
\end{definitionbox}

If such an \(E\) is known, one can attempt to solve

\[
P(D)u=f
\]

by convolution:

\[
\boxed{
u=E*f.
}
\]

Formally,

\[
P(D)(E*f)
=
(P(D)E)*f
=
\delta*f
=
f.
\]

%%%%%%%%%%%%%%%%%%%%%%%%%%%%%%%%%%%%%%%%%%%%%%%%%%%%%%%%%%%%
\subsection{Fundamental Solution of the Derivative}
\label{subsec:fundamental-solution-derivative}
%%%%%%%%%%%%%%%%%%%%%%%%%%%%%%%%%%%%%%%%%%%%%%%%%%%%%%%%%%%%

Since

\[
\boxed{
H'=\delta,
}
\]

the Heaviside distribution is a fundamental solution of the derivative
operator

\[
D=\frac{d}{dx}.
\]

Therefore the equation

\[
u'=f
\]

may be expressed formally as

\[
\boxed{
u=H*f
}
\]

up to the expected homogeneous contribution.

%%%%%%%%%%%%%%%%%%%%%%%%%%%%%%%%%%%%%%%%%%%%%%%%%%%%%%%%%%%%
\subsection{Fundamental Solution of the Laplacian}
\label{subsec:fundamental-solution-laplacian}
%%%%%%%%%%%%%%%%%%%%%%%%%%%%%%%%%%%%%%%%%%%%%%%%%%%%%%%%%%%%

For dimensions

\[
n\geq3,
\]

a fundamental solution for the Laplacian has the form

\[
\boxed{
\Phi(x)
=
\frac{1}{(2-n)\omega_n}
|x|^{2-n},
}
\]

with the convention that \(\omega_n\) denotes the surface area of the
unit sphere in \(\R^n\).

Then

\[
\boxed{
\Delta\Phi=\delta
}
\]

in the sense of distributions.

In two dimensions,

\[
\boxed{
\Phi(x)
=
\frac{1}{2\pi}\log|x|
}
\]

satisfies

\[
\boxed{
\Delta\Phi=\delta.
}
\]

Sign conventions vary depending on whether the operator is taken to be
\(\Delta\) or \(-\Delta\).

%%%%%%%%%%%%%%%%%%%%%%%%%%%%%%%%%%%%%%%%%%%%%%%%%%%%%%%%%%%%
\subsection{Distributional Differential Equations}
\label{subsec:distributional-differential-equations}
%%%%%%%%%%%%%%%%%%%%%%%%%%%%%%%%%%%%%%%%%%%%%%%%%%%%%%%%%%%%

A differential equation

\[
\boxed{
P(D)u=f
}
\]

can be interpreted in \(\mathcal{D}'(\Omega)\).

This means

\[
\boxed{
\langle P(D)u,\varphi\rangle
=
\langle f,\varphi\rangle
}
\]

for every test function \(\varphi\).

This interpretation allows equations to admit solutions with jumps,
corners, point sources, or other singularities.

%%%%%%%%%%%%%%%%%%%%%%%%%%%%%%%%%%%%%%%%%%%%%%%%%%%%%%%%%%%%
\subsection{Point Sources}
\label{subsec:distribution-point-sources}
%%%%%%%%%%%%%%%%%%%%%%%%%%%%%%%%%%%%%%%%%%%%%%%%%%%%%%%%%%%%

A PDE such as

\[
\boxed{
-\Delta u=\delta_0
}
\]

models a source concentrated at one point.

Classically, the right-hand side is not an ordinary function.

Distribution theory gives the equation a precise mathematical meaning.

Such point-source models appear in

\[
\boxed{
\text{electrostatics},
\quad
\text{gravitation},
\quad
\text{heat flow},
\quad
\text{wave propagation}.
}
\]

%%%%%%%%%%%%%%%%%%%%%%%%%%%%%%%%%%%%%%%%%%%%%%%%%%%%%%%%%%%%
\subsection{Green's Functions}
\label{subsec:distribution-greens-functions}
%%%%%%%%%%%%%%%%%%%%%%%%%%%%%%%%%%%%%%%%%%%%%%%%%%%%%%%%%%%%

A Green's function is closely related to a fundamental solution but also
incorporates domain geometry and boundary conditions.

Schematically,

\[
\boxed{
L_xG(x,y)
=
\delta_y(x).
}
\]

If \(G\) is known, solutions may often be represented by integrals
involving

\[
G(x,y).
\]

Thus the Dirac delta serves as the mathematical representation of a unit
point source.

%%%%%%%%%%%%%%%%%%%%%%%%%%%%%%%%%%%%%%%%%%%%%%%%%%%%%%%%%%%%
\subsection{Tensor Products of Distributions}
\label{subsec:tensor-products-distributions}
%%%%%%%%%%%%%%%%%%%%%%%%%%%%%%%%%%%%%%%%%%%%%%%%%%%%%%%%%%%%

If

\[
S\in\mathcal{D}'(\Omega_1)
\]

and

\[
T\in\mathcal{D}'(\Omega_2),
\]

one can define a tensor-product distribution

\[
\boxed{
S\otimes T
}
\]

on

\[
\Omega_1\times\Omega_2.
\]

For regular distributions, this corresponds to products such as

\[
f(x)g(y).
\]

Tensor products allow one-dimensional distributions to generate
higher-dimensional examples.

%%%%%%%%%%%%%%%%%%%%%%%%%%%%%%%%%%%%%%%%%%%%%%%%%%%%%%%%%%%%
\subsection{Schwartz Space}
\label{subsec:schwartz-space}
%%%%%%%%%%%%%%%%%%%%%%%%%%%%%%%%%%%%%%%%%%%%%%%%%%%%%%%%%%%%

Fourier analysis requires a different test-function space.

\begin{definitionbox}{Schwartz Space}
The Schwartz space \(\mathcal{S}(\R^n)\) consists of smooth functions
\(\varphi\) such that

\[
\boxed{
\sup_{x\in\R^n}
|x^\alpha\partial^\beta\varphi(x)|
<
\infty
}
\]

for every pair of multi-indices \(\alpha,\beta\).
\end{definitionbox}

Such functions and all of their derivatives decay faster than every
inverse polynomial.

Typical examples include Gaussian functions.

%%%%%%%%%%%%%%%%%%%%%%%%%%%%%%%%%%%%%%%%%%%%%%%%%%%%%%%%%%%%
\subsection{Tempered Distributions}
\label{subsec:tempered-distributions}
%%%%%%%%%%%%%%%%%%%%%%%%%%%%%%%%%%%%%%%%%%%%%%%%%%%%%%%%%%%%

\begin{definitionbox}{Tempered Distribution}
A tempered distribution is a continuous linear functional on the
Schwartz space:

\[
\boxed{
T:\mathcal{S}(\R^n)\to\mathbb{F}.
}
\]
\end{definitionbox}

The space of tempered distributions is denoted

\[
\boxed{
\mathcal{S}'(\R^n).
}
\]

Tempered distributions are particularly important because the Fourier
transform extends naturally to them.

%%%%%%%%%%%%%%%%%%%%%%%%%%%%%%%%%%%%%%%%%%%%%%%%%%%%%%%%%%%%
\subsection{Examples of Tempered Distributions}
\label{subsec:tempered-examples}
%%%%%%%%%%%%%%%%%%%%%%%%%%%%%%%%%%%%%%%%%%%%%%%%%%%%%%%%%%%%

Important tempered distributions include

\[
\boxed{
\delta,
}
\]

\[
\boxed{
\delta',
}
\]

polynomials,

\[
\boxed{
1,
\quad
x,
\quad
x^2,
\ldots,
}
\]

and many locally integrable functions having at most polynomial growth.

Not every distribution is tempered.

%%%%%%%%%%%%%%%%%%%%%%%%%%%%%%%%%%%%%%%%%%%%%%%%%%%%%%%%%%%%
\subsection{Fourier Transform of Tempered Distributions}
\label{subsec:fourier-tempered-distributions}
%%%%%%%%%%%%%%%%%%%%%%%%%%%%%%%%%%%%%%%%%%%%%%%%%%%%%%%%%%%%

Let \(\mathcal{F}\) denote the Fourier transform using a fixed
normalization convention.

For a tempered distribution \(T\), define its Fourier transform through
duality:

\[
\boxed{
\langle
\mathcal{F}T,\varphi
\rangle
=
\langle
T,\mathcal{F}\varphi
\rangle
}
\]

up to the corresponding normalization convention.

The exact constants and signs depend on the Fourier-transform convention
being used.

The essential point is that the Fourier transform maps

\[
\boxed{
\mathcal{S}'(\R^n)
\to
\mathcal{S}'(\R^n).
}
\]

%%%%%%%%%%%%%%%%%%%%%%%%%%%%%%%%%%%%%%%%%%%%%%%%%%%%%%%%%%%%
\subsection{Fourier Transform of the Delta}
\label{subsec:fourier-delta}
%%%%%%%%%%%%%%%%%%%%%%%%%%%%%%%%%%%%%%%%%%%%%%%%%%%%%%%%%%%%

Under the convention

\[
\widehat{f}(\xi)
=
\int_{\R^n}
f(x)e^{-2\pi i x\cdot\xi}\,dx,
\]

one obtains

\[
\boxed{
\widehat{\delta}=1.
}
\]

Similarly,

\[
\boxed{
\widehat{1}=\delta
}
\]

in the distributional sense.

These identities reveal a duality between complete localization in
physical space and complete delocalization in frequency space.

%%%%%%%%%%%%%%%%%%%%%%%%%%%%%%%%%%%%%%%%%%%%%%%%%%%%%%%%%%%%
\subsection{Fourier Transform and Differentiation}
\label{subsec:fourier-distribution-differentiation}
%%%%%%%%%%%%%%%%%%%%%%%%%%%%%%%%%%%%%%%%%%%%%%%%%%%%%%%%%%%%

Under the same Fourier convention,

\[
\boxed{
\widehat{\partial^\alpha T}(\xi)
=
(2\pi i\xi)^\alpha
\widehat{T}(\xi).
}
\]

Thus differentiation becomes multiplication in frequency space.

Conversely, multiplication by powers of \(x\) corresponds to
differentiation in frequency.

This principle is fundamental to Fourier methods for PDEs.

%%%%%%%%%%%%%%%%%%%%%%%%%%%%%%%%%%%%%%%%%%%%%%%%%%%%%%%%%%%%
\subsection{Solving PDEs by Fourier Transform}
\label{subsec:distribution-fourier-pde}
%%%%%%%%%%%%%%%%%%%%%%%%%%%%%%%%%%%%%%%%%%%%%%%%%%%%%%%%%%%%

Suppose

\[
P(D)u=f.
\]

Taking Fourier transforms formally gives

\[
\boxed{
P(2\pi i\xi)\widehat{u}(\xi)
=
\widehat{f}(\xi).
}
\]

Where division is meaningful,

\[
\boxed{
\widehat{u}(\xi)
=
\frac{\widehat{f}(\xi)}
{P(2\pi i\xi)}.
}
\]

Distribution theory provides a framework for interpreting this formula
when singularities prevent ordinary division or classical solutions.

%%%%%%%%%%%%%%%%%%%%%%%%%%%%%%%%%%%%%%%%%%%%%%%%%%%%%%%%%%%%
\subsection{Principal Value Distribution}
\label{subsec:principal-value-distribution}
%%%%%%%%%%%%%%%%%%%%%%%%%%%%%%%%%%%%%%%%%%%%%%%%%%%%%%%%%%%%

The function

\[
\frac{1}{x}
\]

is not locally integrable across \(x=0\).

Nevertheless, one can define the Cauchy principal-value distribution

\[
\boxed{
\operatorname{p.v.}\frac{1}{x}.
}
\]

Its action is

\[
\boxed{
\left\langle
\operatorname{p.v.}\frac{1}{x},
\varphi
\right\rangle
=
\lim_{\varepsilon\to0^+}
\int_{|x|>\varepsilon}
\frac{\varphi(x)}{x}\,dx.
}
\]

The symmetric exclusion around the singularity produces a meaningful
distribution.

%%%%%%%%%%%%%%%%%%%%%%%%%%%%%%%%%%%%%%%%%%%%%%%%%%%%%%%%%%%%
\subsection{Principal Value and the Hilbert Transform}
\label{subsec:principal-value-hilbert}
%%%%%%%%%%%%%%%%%%%%%%%%%%%%%%%%%%%%%%%%%%%%%%%%%%%%%%%%%%%%

The Hilbert transform involves the singular kernel

\[
\frac{1}{x}.
\]

Formally,

\[
\boxed{
Hf(x)
=
\frac{1}{\pi}
\operatorname{p.v.}
\int_{\R}
\frac{f(y)}
{x-y}\,dy.
}
\]

Distribution theory provides a rigorous interpretation of the singular
kernel.

This connects Distribution Theory with Harmonic Analysis.

%%%%%%%%%%%%%%%%%%%%%%%%%%%%%%%%%%%%%%%%%%%%%%%%%%%%%%%%%%%%
\subsection{Order of a Distribution}
\label{subsec:order-distribution}
%%%%%%%%%%%%%%%%%%%%%%%%%%%%%%%%%%%%%%%%%%%%%%%%%%%%%%%%%%%%

A distribution has finite order locally when its action on test functions
can be bounded using only finitely many derivatives.

Schematically, on a compact set \(K\),

\[
\boxed{
|\langle T,\varphi\rangle|
\leq
C
\sum_{|\alpha|\leq m}
\sup_{x\in K}
|\partial^\alpha\varphi(x)|.
}
\]

The smallest appropriate \(m\), when meaningful in the local setting,
measures the differential complexity of the distribution.

For example, measures define distributions of order zero.

%%%%%%%%%%%%%%%%%%%%%%%%%%%%%%%%%%%%%%%%%%%%%%%%%%%%%%%%%%%%
\subsection{Compactly Supported Distributions}
\label{subsec:compactly-supported-distributions}
%%%%%%%%%%%%%%%%%%%%%%%%%%%%%%%%%%%%%%%%%%%%%%%%%%%%%%%%%%%%

A distribution \(T\) is compactly supported when

\[
\boxed{
\operatorname{supp}(T)
}
\]

is compact.

The space of compactly supported distributions is commonly denoted

\[
\boxed{
\mathcal{E}'(\Omega).
}
\]

Compact support permits distributions to act naturally on arbitrary
smooth functions, not only compactly supported test functions.

%%%%%%%%%%%%%%%%%%%%%%%%%%%%%%%%%%%%%%%%%%%%%%%%%%%%%%%%%%%%
\subsection{Structure of Point-Supported Distributions}
\label{subsec:point-supported-distributions}
%%%%%%%%%%%%%%%%%%%%%%%%%%%%%%%%%%%%%%%%%%%%%%%%%%%%%%%%%%%%

A remarkable theorem states that every distribution supported at a
single point \(a\) is a finite linear combination of derivatives of the
delta distribution.

Thus,

\[
\boxed{
T
=
\sum_{|\alpha|\leq m}
c_\alpha
\partial^\alpha\delta_a.
}
\]

Therefore the delta and its derivatives describe every distribution
concentrated entirely at one point.

%%%%%%%%%%%%%%%%%%%%%%%%%%%%%%%%%%%%%%%%%%%%%%%%%%%%%%%%%%%%
\subsection{Distributional Identities}
\label{subsec:distributional-identities}
%%%%%%%%%%%%%%%%%%%%%%%%%%%%%%%%%%%%%%%%%%%%%%%%%%%%%%%%%%%%

Several useful identities include

\[
\boxed{
H'=\delta,
}
\]

\[
\boxed{
(|x|)''=2\delta,
}
\]

\[
\boxed{
x\delta=0,
}
\]

and

\[
\boxed{
x\delta'=-\delta.
}
\]

For example,

\[
\langle x\delta,\varphi\rangle
=
\langle\delta,x\varphi\rangle
=
0.
\]

Such identities should always be interpreted through their action on
test functions.

%%%%%%%%%%%%%%%%%%%%%%%%%%%%%%%%%%%%%%%%%%%%%%%%%%%%%%%%%%%%
\subsection{Worked Example: Identity \(x\delta'=-\delta\)}
\label{subsec:worked-x-delta-prime}
%%%%%%%%%%%%%%%%%%%%%%%%%%%%%%%%%%%%%%%%%%%%%%%%%%%%%%%%%%%%

\begin{workedexamplebox}{A Distributional Identity}

For any test function \(\varphi\),

\[
\langle x\delta',\varphi\rangle
=
\langle\delta',x\varphi\rangle.
\]

Using

\[
\langle\delta',\psi\rangle
=
-\psi'(0),
\]

we obtain

\[
\langle x\delta',\varphi\rangle
=
-
\left.
\frac{d}{dx}
(x\varphi(x))
\right|_{x=0}.
\]

Since

\[
\frac{d}{dx}
(x\varphi(x))
=
\varphi(x)+x\varphi'(x),
\]

we get

\[
\langle x\delta',\varphi\rangle
=
-\varphi(0).
\]

But

\[
-\varphi(0)
=
\langle-\delta,\varphi\rangle.
\]

Therefore,

\begin{answerbox}
\[
\boxed{
x\delta'=-\delta.
}
\]
\end{answerbox}

\end{workedexamplebox}

%%%%%%%%%%%%%%%%%%%%%%%%%%%%%%%%%%%%%%%%%%%%%%%%%%%%%%%%%%%%
\subsection{Regularization}
\label{subsec:distribution-regularization}
%%%%%%%%%%%%%%%%%%%%%%%%%%%%%%%%%%%%%%%%%%%%%%%%%%%%%%%%%%%%

A distribution can often be approximated by smooth functions.

Let

\[
T\in\mathcal{D}'(\Omega)
\]

and let \(\rho_\varepsilon\) be a mollifier.

Locally one may form

\[
\boxed{
T_\varepsilon
=
T*\rho_\varepsilon.
}
\]

Then \(T_\varepsilon\) is smooth and, on appropriate interior regions,

\[
\boxed{
T_\varepsilon\to T
}
\]

in the distributional sense.

This process is called regularization.

%%%%%%%%%%%%%%%%%%%%%%%%%%%%%%%%%%%%%%%%%%%%%%%%%%%%%%%%%%%%
\subsection{Weak Solutions of Differential Equations}
\label{subsec:distribution-weak-solutions}
%%%%%%%%%%%%%%%%%%%%%%%%%%%%%%%%%%%%%%%%%%%%%%%%%%%%%%%%%%%%

Suppose a classical PDE is

\[
Lu=f.
\]

A weak or distributional solution is an object \(u\) satisfying the
equation after derivatives are transferred to test functions.

For example, if

\[
-u''=f,
\]

then distributionally

\[
\boxed{
-\langle u,\varphi''\rangle
=
\langle f,\varphi\rangle
}
\]

with signs determined by the order of differentiation.

This formulation requires fewer classical derivatives of \(u\).

%%%%%%%%%%%%%%%%%%%%%%%%%%%%%%%%%%%%%%%%%%%%%%%%%%%%%%%%%%%%
\subsection{Distributional Laplacian}
\label{subsec:distributional-laplacian}
%%%%%%%%%%%%%%%%%%%%%%%%%%%%%%%%%%%%%%%%%%%%%%%%%%%%%%%%%%%%

For

\[
T\in\mathcal{D}'(\Omega),
\]

define

\[
\boxed{
\Delta T
=
\sum_{j=1}^{n}
\partial_j^2T.
}
\]

Since second derivatives contribute a positive sign after two
integrations by parts,

\[
\boxed{
\langle\Delta T,\varphi\rangle
=
\langle T,\Delta\varphi\rangle.
}
\]

Thus the Laplacian extends naturally to every distribution.

%%%%%%%%%%%%%%%%%%%%%%%%%%%%%%%%%%%%%%%%%%%%%%%%%%%%%%%%%%%%
\subsection{Distributional Divergence}
\label{subsec:distributional-divergence}
%%%%%%%%%%%%%%%%%%%%%%%%%%%%%%%%%%%%%%%%%%%%%%%%%%%%%%%%%%%%

For a vector-valued locally integrable field

\[
F=(F_1,\ldots,F_n),
\]

the distributional divergence is defined by

\[
\boxed{
\langle
\nabla\cdot F,
\varphi
\rangle
=
-
\int_\Omega
F\cdot\nabla\varphi\,dx.
}
\]

This definition remains meaningful even when the classical derivatives
of \(F\) do not exist.

%%%%%%%%%%%%%%%%%%%%%%%%%%%%%%%%%%%%%%%%%%%%%%%%%%%%%%%%%%%%
\subsection{Distributional Gradient}
\label{subsec:distributional-gradient}
%%%%%%%%%%%%%%%%%%%%%%%%%%%%%%%%%%%%%%%%%%%%%%%%%%%%%%%%%%%%

For a locally integrable scalar function \(u\), the distributional
gradient is

\[
\boxed{
\nabla u
=
(\partial_1u,\ldots,\partial_nu).
}
\]

Each component is defined through

\[
\boxed{
\langle
\partial_ju,\varphi
\rangle
=
-
\int_\Omega
u\,\partial_j\varphi\,dx.
}
\]

This is precisely the derivative notion used in Sobolev spaces.

%%%%%%%%%%%%%%%%%%%%%%%%%%%%%%%%%%%%%%%%%%%%%%%%%%%%%%%%%%%%
\subsection{Distributional Curl}
\label{subsec:distributional-curl}
%%%%%%%%%%%%%%%%%%%%%%%%%%%%%%%%%%%%%%%%%%%%%%%%%%%%%%%%%%%%

In three dimensions, the curl can likewise be defined distributionally.

For a locally integrable vector field \(F\),

\[
\boxed{
\nabla\times F
}
\]

is defined by moving derivatives onto test vector fields through
integration by parts.

Consequently, familiar vector-calculus identities continue to hold
distributionally under suitable interpretations.

%%%%%%%%%%%%%%%%%%%%%%%%%%%%%%%%%%%%%%%%%%%%%%%%%%%%%%%%%%%%
\subsection{Distributional Identities from Vector Calculus}
\label{subsec:distribution-vector-identities}
%%%%%%%%%%%%%%%%%%%%%%%%%%%%%%%%%%%%%%%%%%%%%%%%%%%%%%%%%%%%

Because distributional derivatives commute,

\[
\boxed{
\partial_i\partial_jT
=
\partial_j\partial_iT.
}
\]

Therefore,

\[
\boxed{
\nabla\times(\nabla T)=0,
}
\]

and

\[
\boxed{
\nabla\cdot(\nabla\times F)=0
}
\]

remain valid in the distributional setting.

This is one reason distributions are so effective in PDE theory.

%%%%%%%%%%%%%%%%%%%%%%%%%%%%%%%%%%%%%%%%%%%%%%%%%%%%%%%%%%%%
\subsection{Distribution Theory and Operator Theory}
\label{subsec:distribution-operator-theory}
%%%%%%%%%%%%%%%%%%%%%%%%%%%%%%%%%%%%%%%%%%%%%%%%%%%%%%%%%%%%

Differential operators such as

\[
D,
\qquad
\Delta,
\qquad
-\Delta+V
\]

are often unbounded operators on function spaces.

Distribution theory provides enlarged domains in which these operators
can be interpreted algebraically.

Operator theory then studies properties such as

\[
\boxed{
\text{spectrum},
\quad
\text{self-adjointness},
\quad
\text{resolvents}.
}
\]

Thus the two theories complement each other.

%%%%%%%%%%%%%%%%%%%%%%%%%%%%%%%%%%%%%%%%%%%%%%%%%%%%%%%%%%%%
\subsection{Distribution Theory and Variational Methods}
\label{subsec:distribution-variational-methods}
%%%%%%%%%%%%%%%%%%%%%%%%%%%%%%%%%%%%%%%%%%%%%%%%%%%%%%%%%%%%

Variational problems often produce weak Euler--Lagrange equations.

For example,

\[
J[u]
=
\int_\Omega
\left(
\frac12|\nabla u|^2-fu
\right)\,dx
\]

produces

\[
\boxed{
-\Delta u=f
}
\]

in a weak sense.

Distribution theory provides the language needed to interpret the
derivatives appearing in this equation.

%%%%%%%%%%%%%%%%%%%%%%%%%%%%%%%%%%%%%%%%%%%%%%%%%%%%%%%%%%%%
\subsection{Singular Support and Propagation}
\label{subsec:singular-support-propagation}
%%%%%%%%%%%%%%%%%%%%%%%%%%%%%%%%%%%%%%%%%%%%%%%%%%%%%%%%%%%%

Differential equations often transport or transform singularities.

If \(u\) solves

\[
P(D)u=f,
\]

one may ask how

\[
\operatorname{sing\,supp}(u)
\]

is related to

\[
\operatorname{sing\,supp}(f).
\]

More advanced analysis refines singular support into the
\emph{wavefront set}, which records both

\[
\boxed{
\text{where a singularity occurs}
}
\]

and

\[
\boxed{
\text{in which frequency directions it occurs}.
}
\]

This leads into microlocal analysis.

%%%%%%%%%%%%%%%%%%%%%%%%%%%%%%%%%%%%%%%%%%%%%%%%%%%%%%%%%%%%
\subsection{Common Errors}
\label{subsec:distribution-theory-common-errors}
%%%%%%%%%%%%%%%%%%%%%%%%%%%%%%%%%%%%%%%%%%%%%%%%%%%%%%%%%%%%

\subsubsection{Treating the Dirac Delta as an Ordinary Function}

The expression

\[
\delta(x)
\]

is symbolic.

Its rigorous meaning is

\[
\boxed{
\langle\delta,\varphi\rangle
=
\varphi(0).
}
\]

%%%%%%%%%%%%%%%%%%%%%%%%%%%%%%%%%%%%%%%%%%%%%%%%%%%%%%%%%%%%
\subsubsection{Trying to Evaluate a General Distribution at a Point}

A distribution acts on test functions.

An expression such as

\[
T(x_0)
\]

usually has no meaning for an arbitrary distribution.

%%%%%%%%%%%%%%%%%%%%%%%%%%%%%%%%%%%%%%%%%%%%%%%%%%%%%%%%%%%%
\subsubsection{Forgetting the Minus Sign in the Derivative}

The distributional derivative satisfies

\[
\boxed{
\langle T',\varphi\rangle
=
-\langle T,\varphi'\rangle.
}
\]

For an \(m\)th derivative, the factor is

\[
\boxed{
(-1)^m.
}
\]

%%%%%%%%%%%%%%%%%%%%%%%%%%%%%%%%%%%%%%%%%%%%%%%%%%%%%%%%%%%%
\subsubsection{Assuming Every Distribution Comes from a Function}

Regular distributions arise from locally integrable functions, but
singular distributions such as \(\delta\) do not.

%%%%%%%%%%%%%%%%%%%%%%%%%%%%%%%%%%%%%%%%%%%%%%%%%%%%%%%%%%%%
\subsubsection{Assuming Every Distribution Is Tempered}

Tempered distributions form the smaller space

\[
\boxed{
\mathcal{S}'(\R^n)
}
\]

and satisfy additional growth restrictions.

%%%%%%%%%%%%%%%%%%%%%%%%%%%%%%%%%%%%%%%%%%%%%%%%%%%%%%%%%%%%
\subsubsection{Multiplying Arbitrary Distributions}

Products such as

\[
\delta^2
\]

are not generally defined in classical distribution theory.

Multiplication by a smooth function is well-defined.

%%%%%%%%%%%%%%%%%%%%%%%%%%%%%%%%%%%%%%%%%%%%%%%%%%%%%%%%%%%%
\subsubsection{Confusing Weak Derivatives with Pointwise Derivatives}

A weak derivative is defined through integration against test functions.

It need not exist pointwise everywhere.

%%%%%%%%%%%%%%%%%%%%%%%%%%%%%%%%%%%%%%%%%%%%%%%%%%%%%%%%%%%%
\subsubsection{Ignoring Almost-Everywhere Equivalence}

Locally integrable functions that agree almost everywhere define the same
regular distribution.

%%%%%%%%%%%%%%%%%%%%%%%%%%%%%%%%%%%%%%%%%%%%%%%%%%%%%%%%%%%%
\subsubsection{Ignoring Support Conditions in Convolution}

Convolution of arbitrary distributions is not always defined.

Additional support conditions are required.

%%%%%%%%%%%%%%%%%%%%%%%%%%%%%%%%%%%%%%%%%%%%%%%%%%%%%%%%%%%%
\subsubsection{Confusing Distributional Convergence with Pointwise Convergence}

Distributional convergence means

\[
\boxed{
\langle T_n,\varphi\rangle
\to
\langle T,\varphi\rangle
}
\]

for every test function.

It does not require pointwise convergence.

%%%%%%%%%%%%%%%%%%%%%%%%%%%%%%%%%%%%%%%%%%%%%%%%%%%%%%%%%%%%
\subsubsection{Ignoring Fourier Normalization}

Formulas such as

\[
\widehat{\delta}=1
\]

are structurally stable, but constants in differentiation, inversion,
and convolution formulas depend on the chosen Fourier-transform
convention.

%%%%%%%%%%%%%%%%%%%%%%%%%%%%%%%%%%%%%%%%%%%%%%%%%%%%%%%%%%%%
\subsection{Applications}
\label{subsec:distribution-theory-applications}
%%%%%%%%%%%%%%%%%%%%%%%%%%%%%%%%%%%%%%%%%%%%%%%%%%%%%%%%%%%%

\begin{applicationbox}
\textbf{Partial Differential Equations.}
Distributional solutions allow PDEs to be solved when classical
derivatives do not exist.

\medskip

\textbf{Point Sources.}
Dirac deltas model concentrated sources, loads, charges, masses, and
impulses.

\medskip

\textbf{Fourier Analysis.}
Tempered distributions extend the Fourier transform to deltas,
polynomials, singular kernels, and generalized solutions.

\medskip

\textbf{Harmonic Analysis.}
Principal-value distributions provide rigorous interpretations of
singular integral kernels.

\medskip

\textbf{Potential Theory.}
Fundamental solutions of the Laplacian involve singular distributions and
point sources.

\medskip

\textbf{Operator Theory.}
Distribution spaces provide generalized domains for differential
operators.

\medskip

\textbf{Calculus of Variations.}
Weak Euler--Lagrange equations are naturally interpreted
distributionally.

\medskip

\textbf{Sobolev Spaces.}
Weak derivatives used in Sobolev spaces are distributional derivatives
represented by functions.

\medskip

\textbf{Mechanics.}
Concentrated forces and instantaneous impulses can be represented using
delta distributions.

\medskip

\textbf{Electromagnetism.}
Point charges and singular current distributions naturally produce
distributional source terms.

\medskip

\textbf{Signal Processing.}
The impulse signal is modeled mathematically by the Dirac delta.

\medskip

\textbf{Control Theory.}
Impulse inputs and impulse responses are expressed using distributions.

\medskip

\textbf{Microlocal Analysis.}
Refinements of distribution theory track the location and direction of
singularities.
\end{applicationbox}

%%%%%%%%%%%%%%%%%%%%%%%%%%%%%%%%%%%%%%%%%%%%%%%%%%%%%%%%%%%%
\subsection{Exercises}
\label{subsec:distribution-theory-exercises}
%%%%%%%%%%%%%%%%%%%%%%%%%%%%%%%%%%%%%%%%%%%%%%%%%%%%%%%%%%%%

\begin{exercise}[1]
Explain why

\[
C_c^\infty(\Omega)
\]

is an appropriate space of test functions.
\end{exercise}

\begin{exercise}[1]
For

\[
f(x)=x^2,
\]

write the corresponding regular distribution.
\end{exercise}

\begin{exercise}[1]
Show that two locally integrable functions that agree almost everywhere
define the same distribution.
\end{exercise}

\begin{exercise}[1]
Compute

\[
\langle\delta_2,\varphi\rangle.
\]
\end{exercise}

\begin{exercise}[1]
Compute

\[
\langle\delta',\varphi\rangle.
\]
\end{exercise}

\begin{exercise}[1]
Compute

\[
\langle\delta'',\varphi\rangle.
\]
\end{exercise}

\begin{exercise}[2]
Prove distributionally that

\[
H'=\delta.
\]
\end{exercise}

\begin{exercise}[2]
Show that

\[
(|x|)''=2\delta.
\]
\end{exercise}

\begin{exercise}[2]
Prove that

\[
x\delta=0.
\]
\end{exercise}

\begin{exercise}[2]
Prove that

\[
x\delta'=-\delta.
\]
\end{exercise}

\begin{exercise}[2]
Show that

\[
x^2\delta'=0.
\]
\end{exercise}

\begin{exercise}[2]
Let \(a\) be smooth. Prove that

\[
a\delta_0
=
a(0)\delta_0.
\]
\end{exercise}

\begin{exercise}[2]
Prove the distributional product rule

\[
D(aT)=a'T+aT'.
\]
\end{exercise}

\begin{exercise}[2]
Show that

\[
\delta(ax)
=
\frac{1}{|a|}\delta(x)
\]

for \(a\neq0\).
\end{exercise}

\begin{exercise}[2]
Show that

\[
\operatorname{supp}(\delta_a)
=
\{a\}.
\]
\end{exercise}

\begin{exercise}[2]
Find the singular support of the Heaviside distribution.
\end{exercise}

\begin{exercise}[2]
Verify that

\[
\rho_\varepsilon\to\delta
\]

in the sense of distributions for a standard mollifier.
\end{exercise}

\begin{exercise}[2]
Show that distributional differentiation is continuous with respect to
distributional convergence.
\end{exercise}

\begin{exercise}[2]
Prove that distributional partial derivatives commute.
\end{exercise}

\begin{exercise}[2]
Show that a classical derivative agrees with the corresponding weak
derivative.
\end{exercise}

\begin{exercise}[2]
Determine whether the Heaviside function belongs locally to
\(W^{1,1}(\R)\).
Explain your answer using its distributional derivative.
\end{exercise}

\begin{exercise}[3]
Suppose \(f\) is piecewise \(C^1\) with a jump at \(a\). Derive the
formula

\[
Df
=
f'_{\mathrm{classical}}
+
[f]_a\delta_a.
\]
\end{exercise}

\begin{exercise}[3]
Prove that convolution with a test function smooths a distribution.
\end{exercise}

\begin{exercise}[3]
Show under suitable hypotheses that

\[
\partial^\alpha(T*\varphi)
=
T*(\partial^\alpha\varphi).
\]
\end{exercise}

\begin{exercise}[3]
Show that the Heaviside distribution is a fundamental solution for
\(D=d/dx\).
\end{exercise}

\begin{exercise}[3]
Verify distributionally that

\[
\Delta
\left(
\frac{1}{2\pi}\log|x|
\right)
=
\delta
\]

in \(\R^2\), using an appropriate limiting argument.
\end{exercise}

\begin{exercise}[3]
Derive the weak formulation of

\[
-\Delta u=f.
\]
\end{exercise}

\begin{exercise}[3]
Under the Fourier convention used in this section, prove

\[
\widehat{\delta}=1.
\]
\end{exercise}

\begin{exercise}[3]
Show distributionally that

\[
\widehat{\partial_jT}
=
2\pi i\xi_j\widehat{T}.
\]
\end{exercise}

\begin{exercise}[3]
Verify that the principal-value expression

\[
\operatorname{p.v.}\frac{1}{x}
\]

defines a distribution.
\end{exercise}

\begin{exercise}[3]
Show that every finite measure defines a distribution of order zero.
\end{exercise}

\begin{exercise}[4]
Prove that every distribution has distributional derivatives of all
orders.
\end{exercise}

\begin{exercise}[4]
Prove that every distribution supported at a single point is a finite
linear combination of derivatives of the delta distribution.
\end{exercise}

\begin{exercise}[4]
Study the construction of fundamental solutions for
constant-coefficient differential operators.
\end{exercise}

\begin{exercise}[4]
Use Fourier transforms of tempered distributions to construct a
fundamental solution for a differential operator.
\end{exercise}

\begin{exercise}[4]
Study the relationship between singular support and differential
operators.
\end{exercise}

\begin{exercise}[4]
Investigate wavefront sets as a refinement of singular support.
\end{exercise}

\begin{exercise}[4]
Study the relationship between weak solutions, distributional solutions,
and Sobolev solutions for elliptic PDEs.
\end{exercise}

%%%%%%%%%%%%%%%%%%%%%%%%%%%%%%%%%%%%%%%%%%%%%%%%%%%%%%%%%%%%
\subsection{Conceptual Summary}
\label{subsec:distribution-theory-summary}
%%%%%%%%%%%%%%%%%%%%%%%%%%%%%%%%%%%%%%%%%%%%%%%%%%%%%%%%%%%%

Distribution theory extends the concept of a function by studying linear
actions on smooth test functions.

The test-function space is

\[
\boxed{
\mathcal{D}(\Omega)
=
C_c^\infty(\Omega).
}
\]

A distribution is a continuous linear functional

\[
\boxed{
T:\mathcal{D}(\Omega)\to\mathbb{F}.
}
\]

Its action is written

\[
\boxed{
\langle T,\varphi\rangle.
}
\]

A locally integrable function defines a regular distribution through

\[
\boxed{
\langle T_f,\varphi\rangle
=
\int_\Omega f\varphi\,dx.
}
\]

The Dirac delta is defined by

\[
\boxed{
\langle\delta_a,\varphi\rangle
=
\varphi(a).
}
\]

Every distribution can be differentiated:

\[
\boxed{
\langle
\partial^\alpha T,\varphi
\rangle
=
(-1)^{|\alpha|}
\langle
T,\partial^\alpha\varphi
\rangle.
}
\]

Consequently,

\[
\boxed{
H'=\delta.
}
\]

Weak derivatives are distributional derivatives that happen to be
represented by locally integrable functions.

Convergence in distributions means

\[
\boxed{
\langle T_n,\varphi\rangle
\to
\langle T,\varphi\rangle
}
\]

for every test function.

Fundamental solutions satisfy

\[
\boxed{
P(D)E=\delta
}
\]

and allow differential equations to be approached through convolution.

The Schwartz space

\[
\boxed{
\mathcal{S}(\R^n)
}
\]

leads to tempered distributions

\[
\boxed{
\mathcal{S}'(\R^n),
}
\]

on which the Fourier transform is naturally defined.

Distribution theory therefore supplies a unified language for

\[
\boxed{
\text{generalized functions},
\quad
\text{weak derivatives},
\quad
\text{singular sources},
\quad
\text{Fourier analysis},
\quad
\text{PDEs}.
}
\]

%%%%%%%%%%%%%%%%%%%%%%%%%%%%%%%%%%%%%%%%%%%%%%%%%%%%%%%%%%%%
\subsection{Section Connections}
\label{subsec:distribution-theory-connections}
%%%%%%%%%%%%%%%%%%%%%%%%%%%%%%%%%%%%%%%%%%%%%%%%%%%%%%%%%%%%

\chapterconnection{
Distribution Theory extends differentiation beyond classical functions
and supplies the generalized-function framework required throughout
modern analysis. Measure Theory provides regular distributions through
locally finite measures and locally integrable functions. Functional
Analysis interprets distributions as continuous linear functionals on
topological vector spaces. Fourier Analysis extends naturally to tempered
distributions, allowing deltas, singular kernels, and polynomially growing
functions to be transformed. Harmonic Analysis uses principal-value
distributions and singular integral kernels. Operator Theory uses
distribution spaces as generalized domains for differential operators.
Calculus of Variations produces weak Euler--Lagrange equations whose
derivatives are interpreted distributionally. Potential Theory uses
fundamental solutions and Green's functions for elliptic operators.
Differential Equations and PDEs use distributions to describe weak
solutions, impulses, point sources, and singular forcing. These ideas
also lead toward Sobolev spaces, microlocal analysis, geometric analysis,
and mathematical physics.
}

%%%%%%%%%%%%%%%%%%%%%%%%%%%%%%%%%%%%%%%%%%%%%%%%%%%%%%%%%%%%
% SECTION SOURCE TRACKING
%%%%%%%%%%%%%%%%%%%%%%%%%%%%%%%%%%%%%%%%%%%%%%%%%%%%%%%%%%%%

%------------------------------------------------------------
% 5.14 Distribution Theory
%------------------------------------------------------------
%
% Source basis:
%   - Standard distribution and generalized-function theory.
%   - Standard functional analytic treatment of test-function spaces.
%   - Standard distributional methods in Fourier analysis and PDEs.
%
% Used for:
%   - test functions;
%   - compact support;
%   - test-function spaces;
%   - multi-index notation;
%   - distributions;
%   - regular and singular distributions;
%   - locally integrable functions;
%   - Dirac delta distributions;
%   - measures as distributions;
%   - distributional derivatives;
%   - derivatives of delta;
%   - Heaviside functions;
%   - derivatives across jumps;
%   - weak derivatives;
%   - Sobolev-space connections;
%   - multiplication by smooth functions;
%   - distributional product rules;
%   - translation and scaling;
%   - delta of a function;
%   - support and singular support;
%   - distributional convergence;
%   - approximate identities;
%   - mollifiers;
%   - convolution;
%   - differential operators;
%   - fundamental solutions;
%   - Green's functions;
%   - tensor products;
%   - Schwartz space;
%   - tempered distributions;
%   - Fourier transforms of distributions;
%   - principal-value distributions;
%   - distribution order;
%   - compactly supported distributions;
%   - point-supported distributions;
%   - regularization;
%   - weak and distributional solutions;
%   - distributional vector calculus;
%   - singularity analysis preview.
%
% Notes:
%   - Fourier Analysis develops the transform theory used for tempered
%     distributions in greater depth.
%   - Harmonic Analysis develops principal-value and singular integral
%     operators further.
%   - Potential Theory develops fundamental solutions and Green's
%     functions for elliptic operators.
%   - Functional Analysis supplies the dual-space viewpoint.
%   - Operator Theory supplies the framework for unbounded
%     differential operators.
%   - Calculus of Variations supplies weak formulations that naturally
%     use distributional derivatives.
%   - Geometric Analysis later combines weak derivatives, PDEs,
%     geometry, and variational methods.
%
%%%%%%%%%%%%%%%%%%%%%%%%%%%%%%%%%%%%%%%%%%%%%%%%%%%%%%%%%%%%